% !TeX program = xelatex
% UTF-8. Standalone: xelatex RDS.tex (run from this directory).
% To include in an existing XePersian document, define \RDSInNotebook first.
% Source: Google Drive Notebook/RDS; 19 photographs, sorted by filename.
% Handwritten wording and numbers are retained; uncertain readings are marked.
% pic01.jpg = 20261010_045712517.jpg (Drive ID 1HcMMmvanhphVZ4UnXFBniwQZEwiT-DdL)
% pic02.jpg = 20261010_045720871.jpg (Drive ID 18YbAr6hDmAmeV63fVXOkSR9wdVq11q3B)
% pic03.jpg = 20261010_045726703.jpg (Drive ID 1BzubKFs7X2fFJbadsj_aP9QliOmd8DNK)
% pic04.jpg = 20261010_045735159.jpg (Drive ID 1bvw2663joA25jqPCCrZHhy1GQj7u2841)
% pic05.jpg = 20261010_045742654.jpg (Drive ID 1LbRZ2R8PdTOlnQKw1EJYlf0w-shba-yP)
% pic06.jpg = 20261010_045749021.jpg (Drive ID 1G6oh7w13NZ74nQvNMxvC1IOkCUv-Qd9G)
% pic07.jpg = 20261010_045754400.jpg (Drive ID 15OLtfY9w4O98omoDVZHSWbDTvrg9sbH9)
% pic08.jpg = 20261010_045801627.jpg (Drive ID 1udg-MK4a2qdSaThjsvpWIerONc4YYvv_)
% pic09.jpg = 20261010_045808444.jpg (Drive ID 164mUcp3fgZGiUcBPG5yd3D64M_IPX4vq)
% pic10.jpg = 20261010_045816792.jpg (Drive ID 1KsxR5Kp6gWIbGGRNN_gekTbRrJ7aTAM6)
% pic11.jpg = 20261010_045821856.jpg (Drive ID 14KAl5R-z2A0JQFc7P8s2iqvsLf2-phqv)
% pic12.jpg = 20261010_045828539.jpg (Drive ID 1HgJiSQ26rCIc5XRAsEjWySsXIaC_HX6m)
% pic13.jpg = 20261010_045835154.jpg (Drive ID 1l8toUyvMmJtZACBRYVpddkOby2xCMn6J)
% pic14.jpg = 20261010_045843428.jpg (Drive ID 18PP4_kM-0jwYuiVlgWGVQpU08H0sw51U)
% pic15.jpg = 20261010_045848172.jpg (Drive ID 1yBPhU5cHGysznd1xYo-Yi_U63cSqmb6G)
% pic16.jpg = 20261010_045857659.jpg (Drive ID 1Nymso6WHAsWcE82ipltVQLO55iTgnUZz)
% pic17.jpg = 20261010_045903367.jpg (Drive ID 1DYShQZMxP3JnJ2_fRESaDiHCHtg_4K5o)
% pic18.jpg = 20261010_045913319.jpg (Drive ID 1pI-iZjL1__r5LfHaonLIjws-AxJzUb80)
% pic19.jpg = 20261010_045917960.jpg (Drive ID 1TurikKBeDoK7NVWJalaDiPt_Ya8X2lsT)
\ifdefined\RDSInNotebook
  \let\RDSfinish\relax
\else
  \def\RDSfinish{\end{document}}
  \documentclass[12pt,a4paper]{article}
  \usepackage[margin=22mm]{geometry}
  \usepackage{amsmath,amssymb}
  \usepackage{graphicx}
  \usepackage{xcolor}
  \usepackage{longtable,array}
  \usepackage{enumitem}
  \usepackage{xepersian}
  \IfFontExistsTF{Amiri}{\settextfont{Amiri}}{\settextfont[AutoFakeSlant=0.2]{DejaVu Sans}}
  \setlatintextfont{DejaVu Serif}
  \setlength{\parindent}{0pt}
  \setlength{\parskip}{5pt}
  \setlength{\emergencystretch}{3em}
  \begin{document}
\fi
\providecommand{\RDSunclear}[1]{\textbf{[خوانش نامطمئن: #1]}}
\providecommand{\RDSpage}[2]{\clearpage\section*{تصویر #1}\noindent\lr{\texttt{#2}}\par}
\providecommand{\RDSen}[1]{\lr{#1}}
\providecommand{\RDSmath}[1]{\lr{\(#1\)}}

\begin{center}
{\Large\bfseries سندرم دیسترس تنفسی نوزاد}\par
\lr{65 -- Respiratory Distress Syndrome in the Neonate}\par
\lr{Fanaroff [12]}
\end{center}
\textbf{یادداشت رونویسی:} ترتیب متن مطابق تصاویر است. فلش‌های دست‌نویس به صورت رابطه یا فهرست بازنویسی شده‌اند؛ خوانش‌های نامطمئن در متن مشخص‌اند. برای حفظ کامل منبع، هر ۱۹ تصویر در پیوست آمده است. سربرگ چاپی برگه‌ها \lr{Subject:} و \lr{Date: / /} و نشان چاپی پایین برگه \lr{YASHA} است؛ شمارهٔ دست‌نویس صفحه با شمارهٔ تصویر یکسان است. موضوع دست‌نویس در تصاویر ۲، ۳، ۴، ۶ و ۷–۱۹: \lr{65}؛ در تصویر ۱: \lr{Fanaroff [12]}؛ سربرگ موضوع در تصویر ۵ خالی است.

\RDSpage{۱}{pic01.jpg}
\section*{\lr{RDS}}
نام قبلی بیماری \lr{RDS}، \lr{Hyaline Membrane Disease} است.
\subsection*{\lr{Complications}}
\begin{enumerate}
\item \lr{ICH}
\item \lr{Late onset Sepsis}
\item نکروتایزینگ انتروکولیت
\item \lr{Severe ROP} (رتینوپاتی پره‌مچوریتی)
\item \lr{BPD} (دیسپلازی برونکوپولمونری)
\end{enumerate}
کاملاً نمی‌توان گفت که اختلالات فوق عوارض به علت \lr{RDS}، درمان \lr{RDS} و یا کلاً پره‌مچوریتی رخ داده‌اند.
\subsection*{شیوع بیماری}
\begin{itemize}
\item یک از علل مهم مورتالیتی و موربیدیتی در نوزادان پره‌ترم است.
\item تشخیص قطعی: کمبود سورفاکتانت؛ روش‌ها: بیوشیمیایی، پاتولوژیک انجام می‌شود.
\item بدون شواهد بیوشیمیایی از کمبود سورفاکتانت دشوار است که \lr{RDS} را در نوزادان با وزن تولد (\lr{BW}) شدیداً کم به صورت بالینی تشخیص دهیم.
\item به همین دلیل واژهٔ \lr{Respiratory Insufficiency of Prematurity} برای نوزادان نیازمند اکسیژن یا حمایت تنفسی که نتوان شواهد پاتولوژیک \lr{RDS} به کار برده می‌شود.
\item واژهٔ \lr{Respiratory Instability of Prematurity} برای توصیف نوزاد با \lr{BW} شدیداً کم که نیازمند حمایت تنفسی است ولی فاکتورهای دیگری مثل \RDSunclear{ناپایداری مرکز سیستم تنفس نیز بدان‌ها انتساب می‌شود}.
\end{itemize}

\RDSpage{۲}{pic02.jpg}
یادداشت بالای صفحه: \lr{late preterm} \RDSmath{\to 35\text{--}36}.
\begin{itemize}
\item \lr{RDS} اندکی در نوزادان پسر شایع‌تر است.
\item مهم‌ترین ریسک فاکتور: پره‌مچوریتی و \lr{low BW}.
\item سایر ریسک فاکتورهای مهم: نژاد اروپایی، \lr{late preterm}، زایمان الکتیو در فقدان \lr{labor}.
\item توصیه می‌شود زایمان الکتیو پیش از ۳۹ هفتگی انجام نشود.
\item نوزاد با مادر دیابتیک نیز در ریسک \lr{RDS} قرار دارد.
\item بروز \lr{RDS} با \lr{GA} رابطهٔ معکوس دارد.
\item همهٔ نوزادان با \lr{GA} ۲۲–۲۴، \lr{RDS} خواهند داشت.
\item نوزادان \lr{BW 1251--1500}: ۲۵\% آن‌ها \lr{RDS} خواهند داشت.
\item نوزاد اروپایی با \RDSmath{GA=34}: ۱۰\%؛ با \RDSmath{GA=37}: ۱\%.
\end{itemize}
\subsection*{پاتوفیزیولوژی \lr{RDS}}
\begin{itemize}
\item ریهٔ نوزاد \lr{RDS} در مشاهده سرخ و بی‌هوا است؛ به صورت چرم. در مشاهده شبیه بافت کبد است.
\item در نمای میکروسکوپی، آتلکتازی منتشر دیده می‌شود، به گونه‌ای که تعداد کمی از آلوئول‌ها هوا دارند.
\item راه‌های هوایی انتهایی نمای متورم و متسع دارند (بافت ائوزینوفیلیک فیبرینوز + سلول‌های اپیتلیال، دبری‌های سلولی) (وجه تسمیهٔ بیماری غشای هیالین).
\item در فاز ریکاوری، هایپرپلازی سلول‌های آلوئول دیده می‌شود.
\item سلول تیپ \lr{II}: افزایش فعالیت سورفاکتانت.
\item تولید سورفاکتانت یک پروسهٔ دینامیک است و به عوامل مختلفی ارتباط دارد: \lr{pH}، دما، پرفیوژن.
\item به دنبال سرما، هیپوکسی و اسیدوز: کاهش تولید.
\end{itemize}

\RDSpage{۳}{pic03.jpg}
\begin{itemize}
\item غلظت بالای اکسیژن، ولوم تروما، آتلکتروما به دنبال \lr{MV}: تحریک تولید سایتوکاین‌ها، آسیب اپیتلیوم؛ کاهش تولید.
\item لیک پروتئین‌ها مثل فیبرین به درون فضای انترا‌آلوئولار: غیرفعال شدن سورفاکتانت.
\item کمبود سورفاکتانت $\leftarrow$ کاهش کمپلیانس ریه $\leftarrow$ هایپوونتیلیشن آلوئولار $\leftarrow$ \lr{V/Q} میسمچ $\leftarrow$ \lr{RDS} $\leftarrow$ کاهش اکسیژن‌رسانی $\leftarrow$ اسیدوز لاکتیک.
\item هایپوپرفیوژن ریه ثانویه به وازوکانستریکشن $\leftarrow$ تشدید هایپوکسمی به علت شانت \lr{R\,$\to$\,L} در داکتوس آرتریوزوس، فورامن اووال و ریه.
\item نقش نسبی کمبود سورفاکتانت و هایپوپرفیوژن ریه در تظاهر بالینی \lr{RDS} در هر بیمار متفاوت است.
\end{itemize}
\subsection*{وراثت‌پذیری \lr{RDS}}
\begin{itemize}
\item احتمال ابتلای همزمان در دوقلوهای مونوزیگوت از دی‌زیگوت بیشتر است.
\item علل نادر ژنتیکی برای \lr{RDS}: نقص در ژن‌های کدکنندهٔ:
\begin{enumerate}
\item سورفاکتانت پروتئین \lr{B (SFTPB) (SP-B)}.
\item \lr{SP-C}.
\item \lr{ABCA3}: \lr{ATP binding cassette subfamily A}.
\item \lr{TTF-1}: \lr{Thyroid Transcription Factor 1}.
\end{enumerate}
\item \lr{SP-B deficiency}: وراثت مغلوب (\lr{AR}).
\item دیسترس تنفسی شدید بلافاصله پس از تولد.
\item به مداخلات استاندارد پاسخ نمی‌دهد.
\item ۱ مورد در هر ۱–۵ میلیون تولد زنده.
\end{itemize}

\RDSpage{۴}{pic04.jpg}
\begin{itemize}
\item تولید پاتوژنیک واریانت‌های \lr{SFTPC} به صورت غالب (\lr{AD}) در حدود ۰٫۱\% جمعیت دیده می‌شود.
\item به طور شایع باعث \lr{ILD} در شیرخوار (کمتر از ۱۲ ماه) می‌شود.
\item می‌تواند یک سندرم \lr{RDS-like} در نوزادان ایجاد کند.
\item تجمع \lr{proSP-C}های به شکل \RDSunclear{ناحل محلول؛ واژهٔ پیش از «محلول» روشن نیست}: تحریک آپوپتوز و اتوفاژی.
\item فرم \lr{AR} (بای‌اللیک) واریانت‌های پاتوژنیک \lr{ABCA3}: نارسایی تنفسی نوزادی پیش‌رونده، یا عدم کفایت تنفسی مزمن و \lr{ILD} در بالای ۱ ماه.
\item شایع‌ترین فرم آن (واریانت \lr{p.E292V}): ۱ در هر \lr{90,000} تولد زنده.
\item \lr{ABCA3 def} شایع‌ترین اختلال ژنتیکی هموستاز سورفاکتانت است.
\end{itemize}
\subsection*{\lr{Brain--thyroid--lung Syndrome}}
تریاد:
\begin{enumerate}
\item بیماری نورولوژیک: هایپوتونی + تأخیر تکاملی + اختلالات حرکتی.
\item هایپوتیروئیدی همراه با تولد.
\item \lr{RDS}.
\end{enumerate}
\begin{itemize}
\item حتی درگیری یکی از ارگان‌ها نیز می‌تواند نشان‌دهندهٔ شروع بیماری باشد.
\item ممکن است که تنها فقط یکی از ارگان‌ها درگیر باشد.
\item واریانت پاتوژنیک و \lr{AD} از \lr{NKX2-1}.
\item شیوع بیماری نامعلوم است.
\item اختلال در تکامل \RDSunclear{واژهٔ زیر نشان چاپی: «ساختاری»؟} ریه $\pm$ کاهش بیان ژن‌ها مرتبط با سورفاکتانت.
\end{itemize}

\RDSpage{۵}{pic05.jpg}
\subsection*{پیشگیری: تسریع دارویی بلوغ ریه}
\begin{itemize}
\item داروهای مختلفی بررسی شده‌اند (هورمون تیروئید، آمینوفیلین، کلرول‌آمین \RDSunclear{خوانش نام داروی سوم؛ احتمالاً کاتکول‌آمین})؛ \lr{CS} بهترین گزینه برای تسریع بلوغ ریهٔ جنین است.
\item دوز \lr{CS}: \lr{Physiologic Stress Level} $\leftarrow$ القای بلوغ ریه.
\item \lr{CS} باید ۲۴–۴۸ ساعت قبل از زایمان داده شود.
\item حداقل ۲۴ ساعت قبل؛ حداکثر ۷ روز قبل از زایمان؛ بیشترین تأثیر \lr{CS} در \RDSmath{GA<37} است.
\item با توجه به اینکه \lr{CS} حتی اگر زیر ۲۴ ساعت هم تجویز شود باز هم می‌تواند خطر مرگ، \lr{RDS} و \lr{IVH} را کم کند و سودش هم عالی و بی‌خطر است، ما \lr{CS} را به همه می‌دهیم حتی \lr{Immediate Delivery}.
\item استفاده از \lr{CS} برای \lr{late preterm (34--36)} جای بحث دارد.
\item تأثیر کوتاه‌مدت در کاهش عوارض تنفسی دارد ولی تأثیر بلندمدت در تکامل سیستم عصبی دارد که ممکن است مطلوب نباشد.
\item تأثیر بلندمدت و کوتاه‌مدت دوزهای مکرر \lr{CS} در نوزاد پره‌ترم کاملاً مشخص نیست.
\item امروزه باید دائم عوارض بلندمدت روی رشد و تکامل عصبی خیلی جدی سنجید.
\item تأثیر کورس نجات‌دهنده (\lr{Rescue Course}) \lr{C/S} در \lr{PROM} نامعلوم است.
\end{itemize}

\RDSpage{۶}{pic06.jpg}
\begin{itemize}
\item نوزاد پره‌ترم حداقل ۲۳ هفته از \lr{antenatal CS} سود می‌برد.
\item این تأثیر برای زیر ۲۳ هفته قطعی نیست.
\item دادن \lr{antenatal TSH} هیچ کمکی به بلوغ ریه نمی‌کند.
\end{itemize}
\subsection*{تظاهرات بالینی \lr{RDS}}
یافته‌های بالینی کلاسیک:
\begin{enumerate}
\item تنفس‌های گرانتینگ.
\item رتراکشن.
\item نازال فلرینگ.
\item سیانوز.
\item افزایش نیاز به اکسیژن.
\item \lr{CXR} تشخیصی برای \lr{RDS}.
\end{enumerate}
\begin{itemize}
\item علائم با فاصلهٔ کوتاهی از تولد ایجاد می‌شوند.
\item \lr{RR} معمولاً منظم، بالای \lr{60/min} است.
\item نوزاد معمولاً شواهد پیشرفت نارسایی تنفسی پیدا می‌کند و نیازمند اکسیژن می‌شود.
\item وجود اپیزودهای آپنه در مراحل اولیهٔ \lr{RDS}: یک علامت شوم (\lr{Ominous}) است.
\item مطرح‌کنندهٔ ناپایداری دمای بدن، سپسیس، هایپوکسمی، نارسایی تنفسی.
\item در نوزادی که ما زودهنگام از سورفاکتانت و \lr{MV} استفاده کنیم تظاهر و سیر بیماری متفاوت خواهد بود.
\item رتراکشن‌ها به علت نرمی دنده‌ها بسیار پرومیننت هستند در \lr{RDS}؛ کلاً نرم.
\end{itemize}

\RDSpage{۷}{pic07.jpg}
\begin{itemize}
\item گرانتینگ فقط در اوایل است و ممکن است رفع شود؛ کاراکتریستیک \lr{RDS}.
\item به علت بسته‌شدن ناکامل گلوت طی بازدم؛ برای به دام انداختن هوای آلوئولار و حفظ \lr{FRC}: \lr{Functional Residual Capacity}.
\item به علل دیگر هم رخ می‌دهد: هایپوترمی، هایپوگلایسمی، آپنه، بیماری قلبی، سپسیس.
\item سیانوز در \lr{RDS} در زمینهٔ \lr{V/Q mismatch} رخ می‌دهد و با اکسیژنتراپی رفع می‌شود.
\item تلاش تنفسی زیاد که با ونتیلاتور \lr{Async} باشد می‌تواند با خودش کاهش \lr{Output} قلب، پرفیوژن محیطی را $\downarrow$ و سیانوز ایجاد کند.
\item آکروسیانوز دست و پا در نوزاد طبیعی است و نباید با سیانوز مرکزی که همیشه پاتولوژیک است اشتباه گرفته شود.
\item در نوزاد صداهای تنفسی به صورت گسترده در توراکس پخش می‌شوند و برای تشخیص پاتولوژی قابل اعتماد نیستند.
\item هواگیری غیرقرینهٔ بخش‌های مختلف ریه، همچنین افزایش \RDSunclear{اینفلونسی یا اکوردنسی} فضای اینتراتوراسیک می‌تواند منجر به \lr{leak} هوا و پنوموتوراکس شود.
\item کاهش یک‌طرفهٔ صداهای ریوی + شیفت مدیاستن به سمت مقابل، یا کاهش دوطرفهٔ ورود هوا و صداهای ریوی: پنوموتوراکس.
\item شک به پنوموتوراکس در نوزاد: \lr{Immediate Transillumination}.
\item \lr{CXR} در نوزاد: تأیید مکان \lr{ETT}؛ اگر صداهای ریه غیرقرینه باشند.
\end{itemize}

\RDSpage{۸}{pic08.jpg}
\begin{itemize}
\item سوفل \lr{PDA} طی فاز ریکاوری \lr{RDS} قابل سمع است.
\item یعنی زمانی که مقاومت عروق ریوی (\lr{PVR}) کمتر از عروق سیستمیک شود و شانت \lr{L\,$\to$\,R} رخ دهد.
\item \lr{Distant muffled heart Sounds}: شک به پنوموپریکارد / پریکاردیال افیوژن.
\item دق قفسهٔ سینهٔ نوزاد پره‌ترم هیچ ارزش تشخیصی ندارد.
\item علائم نوزاد در اتاق زایمان یا حداقل ظرف ۴ ساعت شروع می‌شوند.
\item سیر \lr{RDS} غیرعارضه‌دار: شدیدترین علائم در روزهای ۲، ۳؛ شروع ریکاوری بعد از ۷۲ ساعت از شروع علائم.
\item درمان با سورفاکتانت معمولاً سیر بیماری را تسریع می‌کند.
\end{itemize}
\subsection*{چه عواملی ریکاوری را به تأخیر می‌اندازند؟}
\begin{itemize}
\item بیماری آنقدر شدید باشد که به \lr{MV} نیاز باشد.
\item عارضه: \lr{Air leak}.
\item شانت \lr{PDA} قابل ملاحظه.
\item عفونت همزمان.
\item علائم اولیهٔ \lr{BPD}.
\item ریکاوری می‌تواند حتی تا ماه‌ها طول بکشد.
\end{itemize}

\RDSpage{۹}{pic09.jpg}
در اغلب بیماران پره‌ترم، ترنزیشن از فاز ریکاوری \lr{RDS} به بیماری مزمن ریوی به طور نامحسوس و تدریجی رخ می‌دهد.
\subsection*{یافته‌های رادیوگرافیک}
\begin{latin}
\begin{itemize}
\item Reticulogranular infiltrates + Air Bronchograms.
\item Diffuse; classic ground-glass appearance.
\item Normally Symmetric and homogeneous; But can be asymmetric.
\end{itemize}
\end{latin}
\begin{itemize}
\item علت الگوی رتیکولوگرانولر: آتلکتازی آلوئول‌ها؛ علت اصلی. تا حدودی \lr{Pulmonary Edema}.
\item علت \lr{air bronchogram}: مشاهدهٔ برونشیول‌های هوادار در پس‌زمینه‌ای از آلوئول‌های بدون هوا.
\item مشاهدهٔ \lr{air bronchogram} لوکالیزه در \lr{left lower lobe} روی نمای قلب \lr{N/L} است.
\item ولی در \lr{RDS} منتشر هستند، خصوصاً در لوب‌های تحتانی.
\item سایز قلب معمولاً \lr{N/L}، اندکی افزایش‌یافته است.
\item وجود کاردیومگالی مطرح‌کنندهٔ یک نارسایی احتقانی قلب در زمینهٔ یک \lr{hemodynamically significant PDA} می‌باشد.
\item پس از تجویز سورفاکتانت، \lr{CXR} بهبود هوادهی آلوئول‌ها به صورت دوطرفه را نشان خواهد داد.
\item اگر غیرقرینه بود احتمالاً نشان‌دهندهٔ این است که سورفاکتانت وارد یک ریه شده است.
\end{itemize}

\RDSpage{۱۰}{pic10.jpg}
\begin{itemize}
\item هرگز از روی \lr{CXR} نمی‌توان \lr{RDS} غیرعفونی یا آسپتیک را از پنومونی نوزادی (به طور شایع در زمینهٔ \lr{GBS} یا گرم‌منفی‌ها (\lr{E. coli})) تشخیص داد.
\item به همین علت به طور گسترده از \lr{ABT} استفاده می‌شود.
\item نوزاد با \lr{RDS} نسبت به نوزاد با وزن مشابه بدون \lr{RDS}، در \lr{CXR} سایهٔ تیموس بزرگ‌تری دارد.
\item بررسی با اکوکاردیوگرافی در مواردی که شک به \lr{HS-PDA} وجود دارد ضروری است:
\begin{itemize}
\item برای تعیین فشار شریان ریوی.
\item ارزیابی عملکرد قلب.
\item رد بیماری قلبی همراه با تولد (\lr{CHD}).
\item \lr{Obstructed total anomalous pulmonary venous return}.
\end{itemize}
\end{itemize}
\subsection*{درمان \lr{RDS}}
\begin{itemize}
\item \lr{Positive Pressure Ventilation (PPV)}: جزء اصلی درمان \lr{RDS} است که باعث باز ماندن آلوئول‌ها می‌شود.
\item شایع‌ترین اپروچ \lr{MV}:
\end{itemize}
\begin{latin}
Time-cycled, Flow triggered, Volume or pressure-targeted mode with pressure support for spontaneous breathing between mandatory breaths.
\end{latin}
\begin{itemize}
\item \lr{NAVA}: \lr{Neurally Adjusted Ventilatory Assist}؛ با فعالیت‌های الکتریکی دیافراگم تنظیم می‌شود.
\item \lr{CPAP}: تمایل داریم از روش‌های کمتر تهاجمی استفاده کنیم، خصوصاً \lr{n-CPAP}؛ \lr{Nasal}.
\end{itemize}

\RDSpage{۱۱}{pic11.jpg}
\subsection*{مزایای \lr{CPAP}}
\begin{latin}
\begin{itemize}
\item Maintenance of a constant airway opening pressure.
\item Establishment and maintenance of FRC.
\item Reduction of obstructions.
\item Improvement of oxygenation.
\item Reduced volutrauma.
\item Reduced risk of infections.
\end{itemize}
\end{latin}
\subsection*{کدام نوزادان به \lr{CPAP} جواب نمی‌دهند؟}
آپنهٔ ماندگار، بالا نیامدن سیچوریشن، برادی‌کاردی، اسیدوز تنفسی.
\subsection*{\lr{NIPPV}}
\begin{latin}
Noninvasive positive pressure ventilation.
\par nCPAP + intermittent ventilator breaths.
\end{latin}
نقطهٔ منفی \lr{NIPPV} ناتوانی در \lr{Sync} شدن است؛ به همین دلیل \lr{noninvasive-NAVA}.
\subsection*{کافئین}
\begin{itemize}
\item یک درمان شایع برای نوزادان پره‌ترم، خصوصاً زیر ۳۴ هفته.
\item ریسک شکست اکستوبیشن را کم می‌کند.
\item موارد آپنه را کاهش می‌دهد.
\item شدت \lr{BPD} را کاهش می‌دهد.
\item در مصرف طولانی‌مدت ایمن است.
\end{itemize}

\RDSpage{۱۲}{pic12.jpg}
\subsection*{درمان با سورفاکتانت}
انواع: خوکی؛ گاوی.
\subsection*{زمان تجویز}
\begin{itemize}
\item پیشگیری (\lr{Prevention}): در اتاق زایمان.
\item \lr{Early treatment}: ظرف چند ساعت اول.
\item \lr{Rescue}: برای \lr{RDS} تثبیت‌شده.
\end{itemize}
\begin{itemize}
\item به دلیل تجویز آنتی‌ناتال \lr{CS} و \lr{nCPAP} حین دلیوری، در حال حاضر تجویز سلکتیو سورفاکتانت در نوزادانی که نیاز به انتوبه پیدا می‌کنند در حال حاضر آلترناتیو ترجیحی نسبت به تجویز پروفیلاکسی است.
\item همهٔ عوارض \lr{RDS} به دنبال سورفاکتانت کم شده‌اند، بجز نرخ بروز \lr{BPD}.
\item ممکن است به دلیل افزایش نرخ بقای نوزادان بسیار پره‌ترم باشد.
\item ولی دلیل منطقی‌تر این است که \lr{BPD} الزاماً عارضهٔ \lr{RDS} نیست.
\item برخی نوزادان \lr{RDS} به درمان با سورفاکتانت پاسخ قابل انتظار را نمی‌دهند:
\begin{itemize}
\item شدت بیماری اولیه.
\item غیرفعال شدن سورفاکتانت به دلیل التهاب، آسیب اکسیدان‌ها، خونریزی، مایع ادم داخل آلوئول.
\end{itemize}
\item درمان با سورفاکتانت در بیمار با خونریزی ریوی منجر به کاهش نیاز به اکسیژن می‌شود.
\end{itemize}

\RDSpage{۱۳}{pic13.jpg}
\subsection*{انواع سورفاکتانت}
\begin{description}
\item[\lr{Beractant} (سوروانتا):] منشأ گاوی.
\begin{itemize}
\item \lr{25 mg/mL}؛ \lr{11--15.5 mg DPPC}؛ \RDSmath{<1\,\mathrm{mg}} \lr{SP-B + SP-C}.
\item دوز \lr{4 mL/kg}، هر \lr{6 h}.
\item ماکزیموم دوز ۴ نوبت.
\end{itemize}
\item[\lr{Calfactant} (اینفاسورف):] منشأ گاوی.
\begin{itemize}
\item \lr{35 mg/mL}؛ \lr{16 mg DPPC}؛ \lr{0.7 mg SP-B + SP-C}.
\item دوز \lr{3 mL/kg}، \lr{q 6--12 h}.
\item ماکزیموم دوز ۴ نوبت.
\end{itemize}
\item[\lr{Poractant Alfa} (کیوروسورف):] منشأ خوکی.
\begin{itemize}
\item \lr{76 mg/mL}؛ \lr{30 mg DPPC}؛ \lr{1 mg SP-B + SP-C}.
\item \lr{2.5 mL/kg (1st dose), 1.25 mL/kg (2nd, 3rd)}.
\item \lr{q 12 h}؛ ماکزیموم دوز ۳ نوبت.
\end{itemize}
\end{description}
\textbf{یادداشت رونویسی:} در دوزهای \lr{Beractant} و \lr{Calfactant}، \lr{mg} خط خورده و \lr{mL} بالای آن نوشته شده است. مقادیر این بخش عین یادداشت منبع‌اند.
\subsection*{متدهای تجویز سورفاکتانت}
\begin{itemize}
\item تنها متد \lr{FDA approved} از طریق \lr{ETT} است.
\item \lr{INSURE technique (INtubation, SURfactant, Extubate)}.
\item چون که \lr{MV} خطر آسیب ریوی دارد؛ ولی تأثیری در کاهش مرگ‌ومیر یا \lr{BPD} ندارد.
\item با استفاده از \lr{MIST} (مینیمالی اینویزیو سورفاکتانت‌تراپی) یا \lr{LISA} (لس اینویزیو سورفاکتانت ادمینیستریشن)، مزایا: سورفاکتانت و \lr{nCPAP} با هم.
\item خطر \lr{BPD} را در مقایسه با \lr{INSURE} کاهش می‌دهند.
\item \lr{LMA} (لارنژیال ماسک ایروی): هیچ‌گونه بهبودی ندارد.
\end{itemize}

\RDSpage{۱۴}{pic14.jpg}
\begin{itemize}
\item تنها روش غیرتهاجمی دادن سورفاکتانت، فرم \lr{Aerosolized} است.
\item دوز مناسب سورفاکتانت \lr{Nebulized} مشخص نیست.
\item سورفاکتانت بیشترین اثر را دارد اگر ظرف ۱–۲ ساعت از تولد در نوزاد با \lr{RDS} تحت \lr{IMV} تجویز شود.
\item در نوزاد تحت \lr{IMV} در \RDSmath{FiO_2=30\%} و نوزاد \lr{NIMV} ۴۰\% است.
\item روش دادن سورفاکتانت بر اساس تجربهٔ پزشک تعیین می‌شود.
\end{itemize}
\subsection*{سایر کاربردهای سورفاکتانت}
\begin{itemize}
\item \lr{CS} سیستمیک خطر \lr{BPD} را کاهش می‌دهد ولی عوارض نورودولوپمنتال دارد.
\item سورفاکتانت + اسپری بودزوناید: عوارض سیستمیک ندارد.
\end{itemize}
\subsection*{درمان با \lr{NO} استنشاقی (\lr{iNO})}
\begin{itemize}
\item برای درمان نارسایی تنفسی هایپوکسیک در نوزاد \lr{near-term} و ترم تأیید شده است.
\item برای \lr{BPD} و \lr{RDS} با \lr{PHTN} در دست بررسی است.
\item کلاً تأیید نشده.
\item در هایپوپلازی ریوی در زمینهٔ الیگوهیدرآمنیوس، \lr{ROM} طول‌کشیده، \lr{RDS} با \lr{PHTN}، ممکن است باعث بهبود اکسیژناسیون گردد.
\item شک به \lr{PHTN}: انجام اکوکاردیوگرافی در نوزاد پره‌ترم.
\end{itemize}
\subsection*{ارزیابی گاز خون محیطی}
\begin{itemize}
\item به صورت دوره‌ای باید \lr{ABG} اخذ شود.
\item بهتر است از \lr{Indwelling umbilical arterial Catheter} استفاده کنیم.
\end{itemize}

\RDSpage{۱۵}{pic15.jpg}
\begin{itemize}
\item به طور کمتر شایع از شریان‌های محیطی (رادیال، پوستریور تیبیال).
\item نوزاد با \lr{RDS} که به \RDSmath{FiO_2} بالا نیاز دارد، حداقل زمانی که شرایط بالینی‌اش تغییر می‌کند باید \lr{ABG} گرفته شود.
\end{itemize}
\begin{latin}
\begin{itemize}
\item \(PaCO_2\) is maintained in the 45--55 mm Hg range.
\item pH at least 7.25.
\end{itemize}
\end{latin}
\subsection*{مانیتورینگ غیرتهاجمی}
\begin{itemize}
\item پالس اکسی‌متری: به اندام تحتانی وصل می‌شود.
\item \RDSmath{SpO_2} نشانگر تضمین‌کنندهٔ اکسیژن‌رسانی کافی بافتی نیست.
\item تغییرات \RDSmath{SpO_2} تا زمانی که \RDSmath{PaO_2>60} باشد ناچیز است.
\item ممکن است \lr{Occult hyperoxia} رخ دهد.
\item هدف \RDSmath{SpO_2}: \lr{Upper 80 to mid 90\%}.
\item \lr{ROP} با \RDSmath{SpO_2} ۸۵–۸۹\% فراوانی کمتر است.
\item \lr{NEC} با \RDSmath{SpO_2} ۹۰–۹۵\% فراوانی کمتر است.
\item اپیزودهای هایپوکسی ممکن است با اتیولوژی‌های مختلف و پیچیده رخ دهند.
\item احتمال بین \RDSmath{PaCO_2<30} با \lr{periventricular leukomalacia} در نوزاد پره‌ترم نشان‌دهندهٔ نیاز به \RDSmath{CO_2} \lr{monitoring} است.
\end{itemize}
\subsection*{\lr{NIRS}}
\lr{Near-Infrared Spectroscopy}.
\begin{itemize}
\item روش غیرتهاجمی ارزیابی اکسیژناسیون مغز است.
\item هایپوکسی مغزی با \lr{ICH} و کاهش فعالیت الکتریکی همراه است.
\end{itemize}

\RDSpage{۱۶}{pic16.jpg}
\subsection*{نمونه‌گیری شریانی}
\begin{itemize}
\item از شریان آمبیلیکال در شرایط کاملاً استریل.
\item البته شریان‌های آمبیلیکال بر اساس آناتومی یافت می‌شوند.
\item یکی از شریان‌ها کاتتریزه می‌شود.
\item تأیید محل کاتتر با سونوگرافی ضروری است.
\end{itemize}
\subsection*{عوارض کاتتر آمبیلیکال}
\begin{itemize}
\item سیانوز اندام تحتانی یا باتک به علت وازواسپاسم یا ترومبوز.
\item بهتر است کاتتر در \RDSmath{T_6\text{--}T_8} گذاشته شود نه ۳ یا ۴ که نزدیک به فورکیشن آئورت است.
\item آنوریسم و دایسکشن شریان آئورت شکمی؛ تشخیص با سونوگرافی.
\item افزایش فشار خون در صورت ایسکمی شریان رنال.
\item عفونت خون.
\item هرچه کاتتر بیشتر باشد خطر بالاتر.
\end{itemize}
\begin{itemize}
\item قبل از تعبیهٔ خط شریانی: انجام تست \lr{Allen}.
\item شایع‌ترین عارضهٔ \lr{line} شریان محیطی: سیانوز دیستال.
\end{itemize}
\subsection*{درمان اسید–باز}
اسیدوز (تنفسی، متابولیک)، خصوصاً در همراهی با هایپوکسی: وازوکانستریکشن شریان ریوی، \lr{V/Q mismatch} و کاهش برون‌ده قلبی.

\RDSpage{۱۷}{pic17.jpg}
\subsection*{اسیدوز متابولیک}
\begin{itemize}
\item هایپوپرفیوژن بافتی، سپسیس، \lr{IVH}.
\item تجویز سدیم بی‌کربنات به صورت روتین هیچ تأثیری ندارد.
\item اصلاح اسیدوز وابسته به علت زمینه‌ای آن است:
\begin{enumerate}
\item اسیدوز لاکتیک.
\item \lr{RTA}.
\end{enumerate}
\end{itemize}
\subsection*{اسیدوز تنفسی}
\begin{itemize}
\item اغلب نیازمند \lr{MV} است.
\item باید \lr{ABG} سریال اخذ گردد.
\item تجویز سدیم بی‌کربنات در اسیدوز تنفسی کنترااندیکاسیون است؛ منجر به \RDSmath{PCO_2\uparrow}.
\end{itemize}
\subsection*{نتیجهٔ قلبی عروقی}
\begin{itemize}
\item \lr{Overdistention} ریه‌ها در \lr{MV}: کاهش بازگشت وریدی $\leftarrow$ کاهش برون‌ده قلبی.
\item در \lr{CXR}: فلت شدن دیافراگم؛ \lr{Squeezed-heart silhouette}.
\item معیارهای کلینیکال برون‌ده قلبی: علائم حیاتی، نبض پریفرال، کپیلاری ریفیل، برون‌ده ادراری.
\item اختلال عملکرد میوکارد به دنبال ایسکمی یا هایپوپرفیوژن: $\downarrow$ پرفیوژن محیطی + اسیدوز لاکتیک + اسیدوز متابولیک.
\end{itemize}

\RDSpage{۱۸}{pic18.jpg}
\begin{itemize}
\item هایپوتنشن (\RDSmath{BP\downarrow}) سیستمیک در مراحل ابتدایی \lr{RDS} شایع است.
\item در \lr{RDS} متوسط تا شدید:
\end{itemize}
\begin{latin}
Continuous Monitoring of systemic BP by arterial Cath.
\end{latin}
\begin{itemize}
\item چک غیرتهاجمی \lr{BP} صرفاً برای موارد \lr{mild}.
\item اغلب نیازمند سرم کریستالوئید و وازوپرسور خواهیم شد.
\item مواردی که \RDSmath{BP\downarrow} علی‌رغم مایع‌درمانی و وازوپرسور (دوپامین) وجود دارد: چک سطح سرمی کورتیزول؛ درمان با استرس‌دوز هیدروکورتیزون.
\item بسته نشدن \lr{PDA} یک عارضهٔ شایع در نوزاد پره‌ترم با \lr{RDS} است؛ با دیسترس تنفسی و اسیدوز همراهی دارد.
\end{itemize}
\subsection*{شواهد بالینی \lr{HS-PDA}}
\begin{itemize}
\item واید شدن پالس پرشر.
\item \lr{Active precordium}.
\item \lr{bounding peripheral pulses}.
\item سوفل در ناحیهٔ ساب‌کلاویکولار چپ و پشت شانه.
\item در \lr{CXR}: کاردیومگالی.
\item در صورت شک بالینی: اکوکاردیوگرام.
\item تصمیم برای بستن دارویی \lr{PDA} با مهارکنندهٔ سنتز پروستاگلاندین (ایندومتاسین یا ایبوپروفن) همچنان قطعی نیست و محل بحث است.
\item بجز دارویی: محدودیت مایع، وازوپرسور، لوپ دیورتیک (فوروزماید).
\end{itemize}

\RDSpage{۱۹}{pic19.jpg}
\subsection*{درمان آنتی‌بیوتیکی}
\begin{itemize}
\item \lr{RDS} در بالین نه در تصویربرداری به طور قطع قابل افتراق از پنومونی نمی‌باشد.
\item شایع‌ترین ارگانیسم‌ها در پنومونی: \lr{GBS}؛ گرم‌منفی‌ها (\lr{E. coli}).
\item نوزاد با دیسترس تنفسی شدید: \lr{ABT} امپریکال.
\item نوزادان ترم و \lr{late preterm} در ریسک \lr{early onset sepsis} (بر اساس \lr{Sepsis risk calculator}): زودتر \lr{ABT} شروع؟
\item نوزاد زیر ۳۴ هفته با خطر پایین برای \lr{early onset sepsis} (\lr{C/S} بدون \lr{labor} و \lr{ROM}، بدون نگرانی برای عفونت انترا‌یوتِرین (\lr{IUI}) و وضعیت بحرانی نوزاد): \lr{ABT} نمی‌دهیم.
\end{itemize}
\subsection*{شواهد عفونت سیستمیک در نوزاد}
\begin{itemize}
\item هایپوتنشن مقاوم.
\item آپنهٔ \lr{early onset}.
\item نوتروفیلی با شیفت به چپ.
\item اسیدوز متابولیک.
\item نوتروپنی.
\end{itemize}
نبود موارد بالا ردکنندهٔ پنومونی یا سپسیس نیست.
\begin{itemize}
\item شایع‌ترین رژیم \lr{ABT}: آمپی‌سیلین + یک آمینوگلیکوزید.
\item قطع بعد از ۳۴ ساعت بر اساس نتیجهٔ کشت خون و سیر بالینی. \RDSunclear{عدد در تصویر «۳۴» خوانده می‌شود؛ احتمال خوانش دیگر «۳۶» وجود دارد.}
\item \lr{ABT} طولانی‌مدت در نوزاد خطر \lr{NEC} را بالا می‌برد.
\end{itemize}

\clearpage
\section*{پیوست: تصاویر کامل منبع}
تصاویر زیر تمام متن دست‌نویس، حاشیه‌ها، اصلاحات، سربرگ و شماره‌های اصلی را حفظ می‌کنند.
% Works when compiled from this directory or from the repository root.
\IfFileExists{images/pic01.jpg}{\def\RDSimagebase{images/}}{%
  \def\RDSimagebase{pediatrics/ped-neonatology/RDS/images/}}
\newcommand{\RDSoriginal}[2]{%
  \clearpage\subsection*{اصل تصویر #1: \lr{\texttt{pic#2.jpg}}}%
  \begin{center}%
  \includegraphics[width=\linewidth,height=0.87\textheight,keepaspectratio]{\RDSimagebase pic#2.jpg}%
  \end{center}}
\RDSoriginal{۱}{01}
\RDSoriginal{۲}{02}
\RDSoriginal{۳}{03}
\RDSoriginal{۴}{04}
\RDSoriginal{۵}{05}
\RDSoriginal{۶}{06}
\RDSoriginal{۷}{07}
\RDSoriginal{۸}{08}
\RDSoriginal{۹}{09}
\RDSoriginal{۱۰}{10}
\RDSoriginal{۱۱}{11}
\RDSoriginal{۱۲}{12}
\RDSoriginal{۱۳}{13}
\RDSoriginal{۱۴}{14}
\RDSoriginal{۱۵}{15}
\RDSoriginal{۱۶}{16}
\RDSoriginal{۱۷}{17}
\RDSoriginal{۱۸}{18}
\RDSoriginal{۱۹}{19}
\RDSfinish
